% ---------------------------------------------------------------------
\section{Delivery and administration: the analysis was there, the conclusion is new}

\Q Should the dosing schedule, and not only the sampling schedule, be designed for
identification?

\D Every case was priced on 20 Sept and stands unchanged on 27 Sept: amplitude alternation,
a longer and a shorter interval, continuous delivery at the same rate, a $D$ sweep from
$\rn{lnew.dsweep_range_lo}$ to $\rn{lnew.dsweep_range_hi}$, and \rn{lnew.dosesens_n}
readings of the dose calibration in Table~\rn{new.tab.dosesens}. New this week: the exposure sweep to the continuous limit
(\texttt{\rn{gen.therapy}}, on main \rn{git.main_exposure}) and the three named readings of
the admissibility bound (\rn{git.main_readings}).

\F Every case above the rule stands in both drafts at the same value. The rows below it
are new.
\vspace{-2pt}
\begin{wide}
\begin{tabularx}{\linewidth}{@{}p{3.0cm}Lp{2.45cm}@{}}
\toprule
case & result & status \\
\midrule
$\pm\rn{lnew.alt_amp}\%$ amplitude alternation & $\rn{lnew.alt_gain}\%$ of precision for a
 $\rn{lnew.alt_cut}\%$ cut in the trough & carried, \S\rn{new.sec.route} \\
a \rn{lnew.h28_word}-day interval & $\rn{lnew.h28_gain}\%$ for an $\rn{lnew.h28_cut}\%$ cut
 & carried, \S\rn{new.sec.route} \\
a \rn{lnew.h7_word}-day interval & doubles the trough, costs $\rn{lnew.h7_cost}\%$
 & carried, \S\rn{new.sec.route} \\
continuous infusion & no trough, no excitation: the identification stops existing
 & carried, \S\rn{new.sec.route} \\
booked $D = \rn{calc.D}$ at nominal capacity & $\rn{new.sd_monthly}$ against a sweep best of
 $\rn{lnew.dsweep_nom_best}$: within \rn{lnew.dsweep_within} per cent of optimal
 & carried, \S\rn{new.sec.scopeprec} \\
best $D$ in the disease range & $\rn{lnew.dsweep_dis_best}$ at
 $D\approx\rn{lnew.dsweep_dis_d}$ against $\rn{lnew.dsweep_dis_booked}$ booked: no dose
 opens a window & carried, \S\rn{new.sec.scopeprec} \\
largest admissible $D$ at $\rn{calc.dis_pct}\%$ & $\rn{lnew.dmax_5pct}$, where the booked
 increment drives the enzyme to $\rn{lnew.enzyme_x_normal}$ times normal; the monthly design
 still returns $\rn{lnew.d2672_sd}$ & value carried, reading new \\
\midrule
\textbf{no interior optimum} & marker exposure falls monotonically to the continuous limit,
 $\Delta\to\rn{new.delta_cont}$: by $\rn{new.exp_fall_nominal}\times$ at nominal
 ($\rn{new.exp_h28_nominal}$ to $\rn{new.exp_cont_nominal}$) and
 $\rn{new.exp_fall_disease}\times$ at $\rn{calc.dis_pct}\%$ ($\rn{new.exp_h28_disease}$ to
 $\rn{new.exp_cont_disease}$). Peak marker and admissibility ratio fall with it
 & \textbf{new}, \S\rn{new.sec.route} \\
\textbf{three readings of the bound} & R1 (inferred normal, dose scaled with $\theta$):
 constant $\rn{new.adm_r1}$, blind to a breach below $\rn{new.ceiling_blind_pct}\%$
 residual. R2 (inferred normal, booked dose): $\rn{new.adm_r2_20}$ rising to
 $\rn{new.adm_r2_10}$ as capacity halves, the one reading that crosses the bound. R3
 (normal fixed at $\rn{calc.theta}$): breaches at every capacity, $\rn{new.adm_r3_100}$ at
 normal. The readings coincide only at $\theta=\rn{calc.theta}$ and
 $\theta=\rn{calc.theta_dis}$; both anchor numbers sit on the latter. The paper reports on
 R2 & \textbf{new}, \S\rn{new.sec.scopeprec} \\
\bottomrule
\end{tabularx}
\end{wide}

\W \emph{\S\rn{new.sec.route}:} new text inside the administration-route paragraph, on the
absence of an interior optimum; the rest, the delivery cases included, is the 20 Sept text. \emph{\S\rn{new.sec.scopeprec}:}
the bound gained the R1/R2/R3 table and the paragraph naming where the readings coincide; the
$D$ sweep and its two values carry over. \emph{Table \rn{new.tab.dosesens}}
(\S\rn{new.sec.admiss}, the dose calibration) is unchanged.

\Out{Optimise the sampling schedule, which is free. Leave the dosing schedule to therapeutic
judgment, and check identifiability before any move to continuous delivery.}
